\documentclass[10pt,a4paper]{amsart}
\usepackage[margin=19mm]{geometry}
\usepackage{amsmath,amssymb,mathtools}
\usepackage[scr=rsfs,cal=euler]{mathalfa}
\usepackage{xfrac}
\usepackage[unicode,hidelinks,bookmarks=false]{hyperref}
\newcommand{\ie}{i.e.\ }
\newcommand{\cf}{cf.\ }
\newcommand{\resp}{respectively\ }
\newcommand{\Subsection}{\subsection}
\providecommand{\nobreakdash}{\nobreak\hbox{-}}
\setlength{\emergencystretch}{2em}
\multlinegap0pt
\usepackage{mathtools}
\usepackage{bm}
\usepackage[all]{xy}
\usepackage{amssymb}
\usepackage{tikz}
\newtheorem{lemma}{Lemma}
\newtheorem{theorem}{Theorem}
\title{Hecke curves in Frobenius strata of moduli space of rank 2 vector bundles}
\author{Lingguang Li, Hongyi Zhang}
\date{Source: arXiv:2602.09879v1 (CC BY 4.0)}
\begin{document}
\maketitle

\noindent\emph{Source and licence.} Hecke curves in Frobenius strata of moduli space of rank 2 vector bundles. Version arXiv:2602.09879v1 (CC BY 4.0), distributed by arXiv under the Creative Commons Attribution 4.0 International licence, http://creativecommons.org/licenses/by/4.0/, as recorded in the arXiv licence metadata for this exact version and shown on the abstract page https://arxiv.org/abs/2602.09879v1. Full licence notice: \texttt{licenses/arxiv-2602.09879-CC-BY-4.0.txt}.\par\smallskip

\noindent\textit{Editorial scope.} Counted unit: the article's theorem theorem (source0.tex lines 197--204). The article's complete original proof of it is delivered with it (source0.tex lines 205--271, one \texttt{proof} environment of 67 lines). No mathematics is written, repaired or re-derived: the statement and the proof are the article's own lines, in the article's own notation, and no retained line is substituted or re-flowed.  Retained uncounted as the source-local context the proof needs: lines 141--143; lines 144--148; lines 151--182; lines 191--204.  Nothing was omitted from any retained span, so the package carries no omission record.  The retained lines cite 3 works, whose entries are transcribed verbatim from the article's own bibliography source in the e-print and delivered with short editorial labels recorded in the item's reference mapping.  No retained line contains a figure environment, an image-inclusion command or drawing code, so no image is delivered and the package declares no asset dependency.  Editorial formatting only, in addition: an AMS \texttt{amsart} preamble that restates the article's own declarations and macros used by the retained lines, namely the environments \texttt{lemma}, \texttt{theorem} on separate counters, together with the standard packages those lines need.  The retained environments print in this document as Lemma 1, Theorem 1 in order of appearance, whereas the article prints the same items with the numbering of its own full document.  The complete native source of the article as distributed in the arXiv e-print for this exact version is bundled unchanged as \texttt{sources/194349/source0.tex}.
\begin{lemma}[Subsection 4.2 in \cite{JRXY06}]\label{lemma3}
    Let $k$ be an algebraically closed field of characteristic $2$, $X$ a smooth projective curve of genus $g\geqslant 2$ over $k$, $L$ a line bundle on $X$ of degree $d$, $V$ a submodule of $F_*L$ with co-length $\ell$, where $1 \leqslant \ell \leqslant g-2$. Then $V$ is stable and $F^*V$ is not semistable.
\end{lemma}

\begin{lemma}[Subsection 4.3 in \cite{JRXY06}]\label{lemma4}
Let $k$ be an algebraically closed field of characteristic $2$, $X$ a smooth projective curve of genus $g\geqslant 2$ over $k$, $E$  a rank 2 degree $d$ stable vector bundle such that $F^*E$ is not semistable. Suppose the Harder-Narasimhan filtration of $F^*E$ is 
$$0= E_0\subset E_1 \subset E_2=F^*E$$
with $\mathrm{deg}(E_1)=d+j$, $1 \leqslant j \leqslant g-1$. Let $L:=F^*E/E_1$ be the quotient line bundle of degree $d-j$. Then, $E$ is a submodule of $F_*L$ with co-length $g-1-j$.
\end{lemma}

\begin{equation*}
    0 \to E \to F_*L \to \textbf{k}(x) \to 0,
\end{equation*}
where $\textbf{k}(x)$ denotes the skyscraper sheaf of degree $1$ supported at $\{x\}$.


\section{Hecke curves}\label{sec3}
Let $k$ be an algebraically closed field of arbitrary characteristic and $X$ be a smooth projective curve of genus $g \geqslant 2$  over $k$. Let $\mathcal{M}^s_X(r,\mathcal{L})$ be the moduli space of stable vector bundles with fixed determinant $\mathcal{L} \in \mathrm{Pic}^d(X)$. 
Since we will see in Section \ref{sec4} that every rational curve in $\mathcal{M}^s_X(r,d)$ is contained in $\mathcal{M}^s_X(r,\mathcal{L})$ for some $\mathcal{L}\in \mathrm{Pic}^d(X)$, we work directly with rational curves in $\mathcal{M}^s_X(r,\mathcal{L})$. A rational curve in $\mathcal{M}^s_X(r,\mathcal{L})$ is a non-trivial morphism $\phi:\mathbb{P}^1 \to \mathcal{M}^s_X(r,\mathcal{L})$. The degree of $\phi$ is defined as $\mathrm{deg}(\phi):=\mathrm{deg}(\phi^*(-K_{\mathcal{M}^s_X(r,\mathcal{L})}))$. In this section, we focus on the study of \emph{Hecke curves}. 


\emph{Hecke transformations} provide a fundamental method for constructing rational curves passing through a fixed point $[E]\in\mathcal{M}^s_X(r,\mathcal{L})$. This section reviews this approach, drawing on results from Section 1 of \cite{Sun05}.
For a vector bundle $E$ over $X$, we denote its fibre at $x\in X$ by $E_x$, and let $\textbf{k}(x)$ denote the skyscraper sheaf of degree $1$ supported on $\{x\}$. For a subspace $W\subset E_x$, there are two types of modifications of $E$, known as \emph{Hecke transformation} (I) and \emph{Hecke transformation} (II) defined as follows:
\begin{itemize}
    \item [(I)] The \emph{Hecke transformation} (I) of $E$ along $W$ at $x \in X$ is defined as $$E^W:=\mathrm{Ker}(E\to (E_x/W)\otimes\textbf{k}(x))$$
    Explicitly, $E^W$ satisfies
    $$0\to E^W \xrightarrow{\phi} E \to (E_x/W)\otimes\textbf{k}(x)\to 0$$
    with $\phi_x(E^W_x)=W \subseteq E_x$.
    \newline
    \item[(II)] For $W^{\bot} \subset E_x^{\vee}$, the subspace annihilated by $W$, let $(E^{\vee})^{W^{\bot}}$ be the Hecke transformation (I) of  $E^\vee$ along $W^{\bot}$, \emph{i.e.} $(E^{\vee})^{W^{\bot}}$ satisfies 
    $$0\to(E^{\vee})^{W^{\bot}} \xrightarrow{\phi} E^{\vee} \to (E_x^{\vee}/W^{\bot})\otimes\textbf{k}(x)\to 0.$$
    Its dual bundle $\widetilde{E^W}:=((E^{\vee})^{W^{\bot}})^{\vee}$ satisfies 
    $$0 \to E \xrightarrow{\psi} \widetilde{E^W} \to (\widetilde{E_{\ x}^W}/(W^{\bot})^{\vee})\otimes \textbf{k}(x) \to 0$$
    with $\mathrm{Ker}(\psi_x: E_x \to \widetilde{E_{\ x}^W})=W$. We call $\widetilde{E^W}$ the  \emph{Hecke transformation} (II) of $E$ along $W$ at $x \in X$.
\end{itemize}
The above procedures actually construct vector bundles $E^W$ and $\widetilde{E^W}$ satisfying
$$E^W \subset E \subset \widetilde{E^W}$$ with degree $\mathrm{deg}(E^W)=\mathrm{deg}(E)-\dim_k(E_x/W)$ and $\mathrm{deg}(\widetilde{E^W})= \mathrm{deg}(E)+\dim_k(E_x/W)$.

For any $[E]\in\mathcal{M}^s_X(r,d)$ and $x \in X$, let $\zeta \subset E_x$ be a 1-dimensional subspace. Hecke transformation (I) produces a vector bundle $E^{\zeta}$ satisfying 
\begin{equation*}
    0 \to E^{\zeta} \to E \to (E_x/\zeta)\otimes \textbf{k}(x) \to 0.
\end{equation*}

\begin{theorem}[Theorem 1 \cite{Sun05}]
If $k=\mathbb{C}$ and $g \geq 3$, then any rational curve $\phi: \mathbb{P}^1 \to \mathcal{M}^s_X(r,\mathcal{L})$ passing through the generic point has degree at least $2r$. It has degree $2r$ if and only if it is a Hecke curve unless $g=3, r=2$ and $d$ is even.
\end{theorem}

\section{Main results}\label{sec4}
The aim of this section is to prove the following theorem: 
\begin{theorem}
Let $k$ be an algebraically closed field of characteristic $2$, and let $X$ be a smooth projective curve of genus $g\geqslant 2$ over $k$. Let $\mathcal{M}^s_X(2,d)$ be the moduli space of stable vector bundles of rank $2$ and degree $d$ over $X$ and $\mathcal{M}_j:=\{ [E]\in \mathcal{M}_X^s(2,d) \mid \mathrm{HNP}(F^*E) \succcurlyeq \mathscr{P}_j \}$, $\mathcal{M}_j(\mathcal{L}):=\mathcal{M}^s_X(2,\mathcal{L}) \cap \mathcal{M}_j$ be the Frobenius strata,  $1\leqslant j \leqslant g-1$. Then
\begin{itemize}
    \item [(1)] There is no rational curve in $\mathcal{M}_{g-1}$;
    \item[(2)] For any rational curve $\phi: \mathbb{P}^1 \to \mathcal{M}^s_X(2,d)$, there exists $\mathcal{L} \in \mathrm{Pic}^d(X)$ such that $\phi(\mathbb{P}^1) \subseteq \mathcal{M}^s_X(2,\mathcal{L}) \subseteq \mathcal{M}^s_X(2,d)$;
    \item[(3)] For any $\mathcal{L}\in \mathrm{Pic}^d(X)$ and any $[E] \in \mathcal{M}_{g-2}(\mathcal{L}) \backslash \mathcal{M}_{g-1}(\mathcal{L})$, there exists a Hecke curve in $\mathcal{M}_{g-2}(\mathcal{L})\backslash \mathcal{M}_{g-1}(\mathcal{L})$ passing through $[E]$.
\end{itemize}
\end{theorem}

\begin{theorem}
Let $k$ be an algebraically closed field of characteristic $2$, and let $X$ be a smooth projective curve of genus $g\geqslant 2$ over $k$. Let $\mathcal{M}^s_X(2,d)$ be the moduli space of stable vector bundles of rank $2$ and degree $d$ over $X$ and $\mathcal{M}_j:=\{ [E]\in \mathcal{M}_X^s(2,d) \mid \mathrm{HNP}(F^*E) \succcurlyeq \mathscr{P}_j \}$, $\mathcal{M}_j(\mathcal{L}):=\mathcal{M}^s_X(2,\mathcal{L}) \cap \mathcal{M}_j$ be the Frobenius strata,  $1\leqslant j \leqslant g-1$. Then
\begin{itemize}
    \item [(1)] There is no rational curve in $\mathcal{M}_{g-1}$;
    \item[(2)] For any rational curve $\phi: \mathbb{P}^1 \to \mathcal{M}^s_X(2,d)$, there exists $\mathcal{L} \in \mathrm{Pic}^d(X)$ such that $\phi(\mathbb{P}^1) \subseteq \mathcal{M}^s_X(2,\mathcal{L}) \subseteq \mathcal{M}^s_X(2,d)$;
    \item[(3)] For any $\mathcal{L}\in \mathrm{Pic}^d(X)$ and any $[E] \in \mathcal{M}_{g-2}(\mathcal{L}) \backslash \mathcal{M}_{g-1}(\mathcal{L})$, there exists a Hecke curve in $\mathcal{M}_{g-2}(\mathcal{L})\backslash \mathcal{M}_{g-1}(\mathcal{L})$ passing through $[E]$.
\end{itemize}
\end{theorem}

\begin{proof}
(1) In the case of characteristic $2$, according to \cite[Theorem 2.5]{Li14}, the isomorphism
\begin{eqnarray}
S^s_{Frob}: \mathrm{Pic}^{d-g+1}(X) & \to & \mathcal{M}_{g-1} \nonumber \\
    \left[L\right] \qquad & \to & \left[F_*L\right] \nonumber
\end{eqnarray}
implies that $\mathcal{M}_{g-1}$ is isomorphic to an abelian variety, which cannot contain any rational curves.

(2) Suppose $\phi: \mathbb{P}^1 \to \mathcal{M}^s_X(2,d)$ is a rational curve. Consider the determinant morphism 
\begin{eqnarray}
    \mathrm{det}: \mathcal{M}^s_X(2,d) &\longrightarrow &  \mathrm{Pic}^d(X) \nonumber \\
  \left[ E \right]  \quad \, &  \longmapsto &   \, \mathrm{det}(E). \nonumber
\end{eqnarray}
The image of $\mathrm{det} \circ \phi$ is constant since $\mathrm{Pic}^d(X)$ is an abelian variety. Thus, any rational curve is contained in $\mathcal{M}^s_X(2,\mathcal{L})$ for some $\mathcal{L} \in \mathrm{Pic}^d(X)$.

(3) For any $[{E}] \in \mathcal{M}_{g-2}(\mathcal{L})\backslash\mathcal{M}_{g-1}(\mathcal{L})$, according to Lemma \ref{lemma4}, there exists $L \in \mathrm{Pic}^{d-g+2}(X)$ such that 
\begin{equation}\label{4.1}
    0 \to E  \xrightarrow{f} F_*L  \to  \textbf{k}(x) \to 0, 
\end{equation}
is exact. Then, $f$ induces an exact sequence of vector spaces  
\begin{equation*}
    E_x  \xrightarrow{f_x} (F_*L)_x \to \textbf{k}(x)_x \to 0,
\end{equation*}
where $\dim_k\textbf{k}(x)_x=1$. 
For $\zeta:=  \ \mathrm{Ker}(f_x) \subset E_x$, Hecke transformation (I) gives a vector bundle $E^{\zeta}$ satisfying 
\begin{equation*}
    0 \to E^{\zeta} \xrightarrow{q} E \to (E_x/\zeta)\otimes \textbf{k}(x) \to 0. 
\end{equation*}
According to the construction, $q_x(E^{\zeta}_x)=\zeta = \mathrm{Ker}(f_x)$ and $g_x:=f_x \circ q_x: E^{\zeta}_x \to (F_*L)_x$
is zero. We denote the unique maximal ideal of $\mathcal{O}_{X,x}$ by $\mathfrak{m}_x$. To avoid confusion, we denote the localization at a point $x\in X$ by $(-)_{\mathfrak{p}_x}$. Locally, the image of 
$g_{\mathfrak{p}_x}:E_{\mathfrak{p}_x}^{\zeta} \to (F_*L)_{\mathfrak{p}_x}$
is contained in $\mathfrak{m}_x \cdot (F_*L)_{\mathfrak{p}_x}$.
Hence, $g:=f\circ q:E^\zeta \to F_*L$ factors through $F_*L(-x) \to F_*L$. The following commutative diagram 
$$\xymatrix{
0 \ar[r] & E^\zeta \ar[r]^{g} \ar[d]^{\cong} & F_*L \ar[r]\ar@{=}[d] & (F_*L)_x \otimes \textbf{k}(x) \ar[r]\ar[r]\ar@{=}[d] & 0\\
0\ar[r]  &F_*L(-x) \ar[r] & F_*L \ar[r] & (F_*L)_x\otimes \textbf{k}(x) \ar[r] & 0}$$
implies $E^\zeta \cong F_*L(-x)$.

For any point $[l] \in \mathrm{Gr}(1,E_x^\zeta)\cong \mathbb{P}^1$, Hecke transformation (II) of $E^\zeta$ along $l\subset E_x^\zeta$ yields a vector bundle $\widetilde{(E^{\zeta})^l}$ satisfying
\begin{equation}\label{applyfunctor}
    0 \to E^\zeta \xrightarrow{r} \widetilde{(E^{\zeta})^l}\xrightarrow{s} (\widetilde{(E^{\zeta})_x^l}/(l^{\bot})^{\vee})\otimes \textbf{k}(x) \to 0.
\end{equation}
Note that for $l = \mathrm{ker}(q_x)$, we have $\widetilde{(E^\zeta)^{l}} \cong E$ by Section \ref{sec3}. 

For an injection of $k$-vector space $i: \widetilde{(E^{\zeta})_x^l}/(l^{\bot})^{\vee} \to (F_*L)_x$, if there exists an injection $h_{\mathfrak{p}_x}$ making the following diagram of $\mathcal{O}_{X,x}$-modules commute, 
$$\xymatrix{
0 \ar[r] & E_{\mathfrak{p}_x}^{\zeta} \cong F_*L(-x)_{\mathfrak{p}_x} \ar[r]^{\qquad r_{\mathfrak{p}_x}} \ar[d]^{\cong} & \widetilde{(E^\zeta)^l}_{\mathfrak{p}_x} \ar[r]^{s_{\mathfrak{p}_x}\quad} \ar@{-->}[d]^{h_{\mathfrak{p}_x}}& \widetilde{(E^{\zeta})_x^l}/(l^{\bot})^{\vee} \ar[r]\ar[d]^{i} & 0\\
0\ar[r]  &F_*L(-x)_{\mathfrak{p}_x} \ar[r]  & (F_*L)_{\mathfrak{p}_x} \ar[r]^{t_{\mathfrak{p}_x}} & (F_*L)_x \ar[r] & 0
}$$
we have an injection 
$$h:\widetilde{(E^\zeta)^l} \to F_*L$$ 
of sheaves of $\mathcal{O}_{X}$-modules and $\mathrm{deg}(F_*L) = \mathrm{deg}(\widetilde{(E^\zeta)^l})+1$.

Since the problem is local, we may assume $(F_*L)_{\mathfrak{p}_x} \cong F_*(k[[t]])$ is a rank 2 free $k[[t^2]]$-module generated by $\{1,t\}$ and $F_*L(-x)_{\mathfrak{p}_x}=\mathfrak{m}_x \cdot (F_*L)_{\mathfrak{p}_x}$ generated by $t^2 \cdot \{1,t\} = \{t^2, t^3\}$. Then $(F_*L)_x$ is a $k$-vector space spanned by $\{\overline{1}, \overline{t} \}$. Denote $\boldsymbol{u}:=r_{\mathfrak{p}_x}(t^2)$ and $\boldsymbol{v}:=r_{\mathfrak{p}_x}(t^3)$.  Suppose $\widetilde{(E^\zeta)^l}_x/(l^{\bot})^{\vee}$ is spanned by  $\overline{\boldsymbol{w}}$ and $\boldsymbol{w}$ is a preimage of $\overline{\boldsymbol{w}}$ in $\widetilde{(E^\zeta)^l}_{\mathfrak{p}_x}$. Then, $\widetilde{(E^\zeta)^l}_{\mathfrak{p}_x}$ can be expressed as 
$$\widetilde{(E^\zeta)^l}_{\mathfrak{p}_x} = k[[t^2]] \cdot \boldsymbol{u}   \oplus k[[t^2]] \cdot \boldsymbol{v}  + k\cdot \boldsymbol{w}.$$
Let $i: \overline{\boldsymbol{w}} \mapsto \overline{a+bt} \in (F_*L)_x$, where $a,b \in k$. Then $a+bt$ is a representative element of $t_{\mathfrak{p}_x}^{-1}(\overline{a+bt})$ in $(F_*L)_{\mathfrak{p}_x}$. We can construct the natural embedding 
\begin{equation*}
    h_{\mathfrak{p}_x}: \boldsymbol{u} \mapsto t^2, \quad  \boldsymbol{v} \mapsto t^3, \quad  \boldsymbol{w} \mapsto a+bt.
\end{equation*}
It is obvious that $h_{\mathfrak{p}_x}$ makes the diagram commute. Since $$\mathrm{deg}(F_*L/\widetilde{(E^\zeta)^l})=1,$$
by Lemma \ref{lemma3}, $\widetilde{(E^\zeta)^l}$ is stable and $F^*(\widetilde{(E^\zeta)^l})$ is not semistable. 

If there exists $l\subset E_x^\zeta$ such that $[\widetilde{(E^\zeta)^l}] \notin \mathcal{M}_{g-2}(\mathcal{L})\backslash\mathcal{M}_{g-1}(\mathcal{L})$, by Lemma \ref{lemma4}, there exists a line bundle $L'$ of $\mathrm{deg}(L')<d-g+2$ such that  $\widetilde{(E^\zeta)^l} \to F_*L'$ is an embedding. But 
$$ \mathrm{deg}(F_*L') = \mathrm{deg}(L') + g-2 \leqslant d-1 < d =\mathrm{deg}(\widetilde{(E^\zeta)^l}),$$
which is a contradiction. Thus, we obtain a Hecke curve $\{\widetilde{(E^\zeta)^l} \mid  [l] \in \mathrm{Gr}(1,E_x^\zeta)\}$ in $\mathcal{M}_{g-2}(\mathcal{L})\backslash\mathcal{M}_{g-1}(\mathcal{L})$ passing through $[E]$.

\end{proof}

\begin{thebibliography}{99}
\bibitem[JRXY06]{JRXY06} K. Joshi and S. Ramanan and E. Z. Xia and J.-K. Yu, {\it On vector bundles destabilized by Frobenius pull-back}, Compos. Math. {\bf 142} (2006), 616--630.
\bibitem[Li14]{Li14} L. Li, {\it The morphism induced by Frobenius push-forward}, Sci. China Math. {\bf 57} (2014), 61-67.
\bibitem[Sun05]{Sun05} X. Sun, {\it Minimal rational curves on moduli spaces of stable bundles}, Math. Ann. {\bf 331} (2005), 925-937.
\end{thebibliography}
\end{document}
